% IEEE Conference Paper Format
% Hate Speech Detection and Highlighting in Vietnamese Social Media Comments
\documentclass[conference]{IEEEtran}
\IEEEoverridecommandlockouts

\usepackage{cite}
\usepackage{amsmath,amssymb,amsfonts}
\usepackage{graphicx}
\usepackage{textcomp}
\usepackage{xcolor}
\usepackage{booktabs}
\usepackage{multirow}
\usepackage{enumitem}
\usepackage{array}
\usepackage{url}
\usepackage[hidelinks]{hyperref}
\usepackage[utf8]{inputenc}
\usepackage[T5,T1]{fontenc}
\usepackage[vietnamese,english]{babel}
\usepackage{tikz}
\usetikzlibrary{shapes.geometric,arrows.meta,positioning,calc,fit}

\def\BibTeX{{\rm B\kern-.05em{\sc i\kern-.025em b}\kern-.08em
    T\kern-.1667em\lower.7ex\hbox{E}\kern-.125emX}}

\title{Hate Speech Detection and Highlighting\\in Vietnamese Social Media Comments}

\author{
\IEEEauthorblockN{An Truong Hoang Thanh\textsuperscript{1,2},~An Nguyen Xuan\textsuperscript{1,2},~Binh Mai Thai\textsuperscript{1,2},~Phat Nguyen Cong\textsuperscript{1,2},~Huong Nguyen Le Quynh\textsuperscript{1,2}}
\IEEEauthorblockA{
  \textsuperscript{1}\textit{University of Information Technology (UIT)} \\
  \textsuperscript{2}\textit{Vietnam National University, Ho Chi Minh City (VNU-HCM)},
  Vietnam \\
  \{23520032, 23520023, 23520158, 23521143, 21520255\}@gm.uit.edu.vn
}
}

\begin{document}
\maketitle

% --------------------------------------------------------------------------
\begin{abstract}
The rise of hate speech on Vietnamese social media has created a pressing need for
automated systems that go beyond simple classification. In this work, we address the
challenge of target-specific and implicit hate speech detection by introducing the HARE
(Hate speech detection with Rationale Extraction) framework for Vietnamese. By
integrating Chain-of-Thought (CoT) reasoning into the Qwen2.5-3B-Instruct model, our
approach moves beyond traditional black-box classifiers: instead of producing only a
label, HARE generates semantic rationales and implied statements that reveal the hidden
meaning behind a comment, via a two-stage QLoRA fine-tuning process on the ViTHSD
dataset. Our experiments show that the proposed model achieves an F1-Micro of
60.26\%, outperforming PhoBERT by +6.14\% and Flan-T5 by +13.42\%, with a
particularly notable gain of +24.87\% on the Politics\#Hate category. We further present
a full-stack real-time deployment architecture built with FastAPI, React, and Apache
Kafka/Spark, equipped with a keyword-highlighting module, demonstrating applicability
to large-scale Vietnamese social media monitoring.
\end{abstract}

\begin{IEEEkeywords}
hate speech detection, Vietnamese NLP, chain-of-thought prompting, QLoRA, explainable AI,
ViTHSD, Qwen2.5
\end{IEEEkeywords}
% --------------------------------------------------------------------------


% ==========================================================================
\section{Introduction}\label{sec:intro}
% ==========================================================================

\subsection{Background and Motivation}

Vietnamese social media platforms---Facebook, YouTube, and TikTok---have become the
primary channels for public discourse, but this openness has simultaneously enabled a
surge in toxic and hateful content. The Vietnamese language presents unique challenges
for automated content moderation: it is monosyllabic and tonal, meaning is
deeply context-dependent, and contemporary online users frequently embed hateful
messages in culturally specific slang (e.g., \foreignlanguage{vietnamese}{\textit{``ba
que,'' ``bò đỏ''}}), teencode, sarcasm, and wordplay that evade keyword-based filters.
Addressing this task is critical: undetected hate speech can incite real-world
discrimination and violence, and manual moderation does not scale.

Existing automated approaches to Vietnamese hate speech detection suffer from three
interconnected limitations. First, \textit{lack of explainability}: deep learning
classifiers operate as black boxes, providing labels without reasoning, which undermines
moderator trust. Second, \textit{inability to detect implicit hate}: most models rely on
surface-level keywords and fail to capture attacks that use cultural metaphors, sarcasm,
and community-specific context rather than explicit profanity. Third, \textit{neglect of
attack target}: prior work focuses on toxicity severity while overlooking \textit{who}
is being attacked---an individual, a group, or a political entity---information essential
for calibrating moderation responses.

\subsection{Task Definition}

We formally define the \textbf{Vietnamese Targeted Hate Speech Detection Task} as
follows. Given a Vietnamese social media comment $c$, the system must predict a set of
labels $\mathcal{L} \subseteq \{T_i\texttt{\#}L_j\}$, where $T_i \in
\{\text{Individuals, Groups, Religion, Race, Politics}\}$ is the targeted entity and
$L_j \in \{\text{Normal, Clean, Offensive, Hate}\}$ is the severity level. This is a
\textit{multi-label} task: a single comment may simultaneously attack multiple targets
at different severity levels. For example, the comment:

\smallskip
\noindent\foreignlanguage{vietnamese}{\textit{``Thằng A ngu thế cũng làm ca sĩ
(\texttt{Individuals\#Offensive}), bọn fan của nó thì cũng súc vật không kém
(\texttt{Groups\#Hate})''}}
\smallskip

\noindent yields two labels: the first clause offends an individual, while the second
expresses hate toward a fan community.

\subsection{Contributions}

To address these limitations, this paper makes five contributions:
\begin{enumerate}[leftmargin=*, label=\textbf{C\arabic*.}]
\item \textbf{Task Systematization}: We formalize the Vietnamese Targeted Hate Speech
Detection task from a target-oriented perspective and provide a comprehensive
linguistic analysis of the ViTHSD dataset, including statistical characteristics, class
imbalance, and the implicit hate challenge.
\item \textbf{HARE Framework Adaptation}: We adapt the HARE framework \cite{Yang2023HARE}
for Vietnamese by engineering a three-phase CoT annotation pipeline using Gemini 2.0
Flash \cite{google2025gemini2flash} to automatically generate high-quality rationales and
implied statements for ambiguous and implicit hate speech examples.
\item \textbf{Two-Stage Semantic Alignment Training}: We propose a two-stage QLoRA
\cite{Dettmers2023QLoRA} fine-tuning strategy for Qwen2.5-3B-Instruct
\cite{Yang2025Qwen25}---Stage~1 learns classification boundaries on 7,540 labeled
samples, while Stage~2 performs semantic alignment on 1,221 rationale-enriched pairs.
\item \textbf{Empirical Evaluation}: We provide a rigorous comparison against four
baselines---PhoBERT \cite{nguyen2020phobert}, Flan-T5 \cite{chung2022scaling}, Qwen2.5
Vanilla (Stage~1 only), and the proposed HARE model---achieving state-of-the-art
results, especially on the low-resource Politics and Groups categories.
\item \textbf{Deployment Architecture}: We present a complete real-time system built on
Apache Kafka, Spark, FastAPI, and React+Vite, enabling scalable Vietnamese social
media monitoring with visual keyword highlighting.
\end{enumerate}

\subsection{Paper Organization}
The remainder of this paper is organized as follows.
Section~\ref{sec:related} reviews related work.
Section~\ref{sec:dataset} analyzes the ViTHSD dataset.
Section~\ref{sec:method} presents the HARE methodology.
Section~\ref{sec:results} reports experimental results.
Section~\ref{sec:deployment} describes the deployment architecture.
Section~\ref{sec:conclusion} concludes the paper and outlines future work.

% ==========================================================================
\section{Related Work}\label{sec:related}
% ==========================================================================

This section reviews the three bodies of work most directly relevant to our approach:
Vietnamese NLP models, Vietnamese hate speech datasets, and explainable hate speech
detection.

\subsection{Vietnamese NLP and Pre-trained Language Models}

Vietnamese NLP has transformed significantly over the past decade. Before 2019, most
systems relied on static word embeddings such as Word2Vec \cite{Mikolov2013Efficient} and
FastText \cite{Bojanowski2017Enriching}. The Transformer architecture
\cite{Vaswani2017Attention} and BERT \cite{Devlin2019BERT} fundamentally changed the
field. PhoBERT \cite{nguyen2020phobert}, trained on 20 GB of formal Vietnamese text,
established new benchmarks on standard NLP tasks including part-of-speech tagging and
named entity recognition. However, PhoBERT exhibits a clear domain gap on informal social
media language: slang, abbreviations, and code-switching common in online discourse were
absent from its pre-training corpus.

Large language models (LLMs) such as Qwen2.5 \cite{Yang2025Qwen25} overcome this gap
through more diverse pre-training. Qwen2.5-3B-Instruct outperforms comparably sized
Llama~2 \cite{llama} and Gemma \cite{Mesnard2024Gemma} on Vietnamese benchmarks, making
it a strong backbone for instruction-following tasks in Vietnamese. This motivates our
choice of Qwen2.5-3B-Instruct over encoder-only alternatives.

\subsection{Hate Speech Datasets for Vietnamese}

ViHSD \cite{10.1007/978-3-030-79457-6_35} was the first large-scale Vietnamese hate
speech dataset, containing over 30,000 comments labeled as Clean, Offensive, or Hate.
While pivotal for establishing baselines, ViHSD provides no information about the
\textit{target} of the hate speech. ViTHSD \cite{vithsd} addresses this limitation by
extending the annotation scheme with five target categories (Individuals, Groups,
Religion, Race/Ethnicity, Politics), following the Aspect-Based Sentiment Analysis
paradigm \cite{Pontiki2014SemEval4} where the model identifies not only the
\textit{polarity} but also the \textit{aspect} being judged. Unlike ViHSD, ViTHSD
supports multi-label annotation, enabling a single comment to carry multiple
target-severity pairs.

\subsection{Explainable AI and Chain-of-Thought Prompting}

Explainability is a critical requirement in content moderation AI. The HARE framework
\cite{Yang2023HARE}, introduced at EMNLP 2023, used LLM-generated rationales to make
hate speech detection interpretable. Chain-of-Thought (CoT) prompting
\cite{Wei2022ChainOfThought} elicits step-by-step reasoning from LLMs, allowing them to
decompose complex semantic tasks into interpretable reasoning chains. Prior work on
English-language hate speech shows that training with CoT rationales significantly
improves performance on implicit and context-dependent examples
\cite{Yang2023HARE,Wei2022ChainOfThought}. Unlike these works, we apply CoT reasoning to
Vietnamese, which requires culturally specific prompt design, Vietnamese-aware data
validation, and semantic alignment tailored to Vietnamese social media pragmatics.

% ==========================================================================
\section{ViTHSD Dataset Analysis}\label{sec:dataset}
% ==========================================================================

This section analyzes the ViTHSD dataset to motivate our modeling choices. We examine
label structure, statistical distribution, class imbalance, and the implicit hate challenge
that drives our rationale augmentation strategy.

\subsection{Dataset Structure and Label Definition}

ViTHSD \cite{vithsd} contains 10,000 Vietnamese social media comments annotated with a
composite label scheme \texttt{[Target]\#[Level]}. Five target categories are defined:
(1) \textbf{Individuals}---attacks targeting a specific person, the most common online
abuse form; (2) \textbf{Groups}---attacks targeting organizations, communities, or fan
groups; (3) \textbf{Religion}---targeting religious beliefs, figures, or communities;
(4) \textbf{Race/Ethnicity}---targeting regional or ethnic identity, the most dangerous
category; (5) \textbf{Politics}---targeting political views, parties, government
policies, or figures. The four severity levels are: \textbf{Normal}~(0)---no negative
content; \textbf{Clean}~(1)---neutral or positive mention of a target;
\textbf{Offensive}~(2)---vulgar or mocking language short of incitement;
\textbf{Hate}~(3)---dehumanizing language that incites discrimination or violence.

\subsection{Statistical Distribution and Data Imbalance}

Table~\ref{tab:target_dist} presents the sample distribution by target category across
train, development, and test splits.

\begin{table}[ht]
\centering
\caption{Distribution of Samples by Target Category}
\label{tab:target_dist}
\scriptsize
\setlength{\tabcolsep}{3pt}
\resizebox{\columnwidth}{!}{%
\begin{tabular}{lrrrrr}
\toprule
\textbf{Target} & \textbf{Train} & \textbf{Dev} & \textbf{Test} & \textbf{Total} & \textbf{\%} \\
\midrule
Individuals  & 5,480 & 938 & 1,398 & 7,816 & $\sim$60\% \\
Groups       & 2,977 & 517 & 769   & 4,263 & $\sim$33\% \\
Race/Ethni.  & 502   & 74  & 129   & 705   & $\sim$5\%  \\
Politics     & 363   & 57  & 89    & 509   & $\sim$4\%  \\
Religion     & 24    & 8   & 6     & 38    & $<$0.3\%   \\
\bottomrule
\end{tabular}
}
\end{table}

The Individuals category dominates at approximately 60\%, reflecting the prevalence of
personal attacks on Vietnamese social media. Politics and Religion are severely
underrepresented (less than 5\% combined). This imbalance is a major training hurdle:
without mitigation, models predict majority-class labels, achieving reasonable accuracy
on common cases but failing almost entirely on rare but high-risk categories. The
Religion class, with only 24 training samples, is nearly intractable without specialized
techniques such as data augmentation or few-shot learning.

Table~\ref{tab:level_dist} shows the distribution by severity level. The Clean category
dominates, requiring models to maintain high Recall to surface Hate labels hidden among
a large volume of Clean comments.

\begin{table}[ht]
\centering
\caption{Distribution of Samples by Severity Level}
\label{tab:level_dist}
\scriptsize
\setlength{\tabcolsep}{3pt}
\resizebox{\columnwidth}{!}{%
\begin{tabular}{lrrrrr}
\toprule
\textbf{Split} & \textbf{Normal} & \textbf{Clean} & \textbf{Offens.} & \textbf{Hate} & \textbf{Total} \\
\midrule
Train & 1,520 & 2,480 & 1,169 & 1,831 & 7,000 \\
Dev   & 263   & 454   & 189   & 295   & 1,201 \\
Test  & 402   & 618   & 324   & 456   & 1,800 \\
\midrule
\textbf{Total} & \textbf{2,185} & \textbf{3,552} & \textbf{1,682} & \textbf{2,582} & \textbf{10,001} \\
\bottomrule
\end{tabular}
}
\end{table}

\subsection{The Implicit Hate Challenge}

A critical characteristic of ViTHSD is the prevalence of \textit{implicit hate
speech}---comments that carry deeply offensive meaning without explicit profanity. Two
representative examples illustrate why semantic reasoning, not keyword matching, is
required:

\begin{enumerate}[leftmargin=*]
\item \foreignlanguage{vietnamese}{\textit{``Optimusthree với 99\% sức mạnh''}}
(``Optimusthree at 99\% power''):
\textit{Surface reading}: a neutral strength statement.
\textit{Cultural context}: a sarcastic gaming meme implying the player is
\textit{extremely weak} (only 1\% capability). Without this cultural knowledge, a
model classifies this as Clean. Correct label: \texttt{Individuals\#Offensive}.

\item \foreignlanguage{vietnamese}{\textit{``Lũ ba que xỏ lá''}}:
``Ba que'' is political slang for supporters of the former South Vietnamese regime
(yellow flag with three red stripes); ``xỏ lá'' means deceitful. Combined, this targets
a political group with hate. Correct label: \texttt{Politics\#Hate}.
\end{enumerate}

These examples motivate enriching training data with \textit{Implied Statements}
(concise paraphrases revealing true intent) and \textit{Rationales} (logical
explanations grounding the classification in text evidence), which we generate in
Section~\ref{sec:method}.








% ==========================================================================
\section{Methodology}\label{sec:method}
% ==========================================================================

The HARE methodology for Vietnamese follows a two-phase pipeline, illustrated in
Fig.~\ref{fig:pipeline}: (1) \textbf{Annotation Phase}---using a commercial LLM to
generate high-quality rationales and implied statements; (2) \textbf{Training
Phase}---fine-tuning Qwen2.5-3B-Instruct in two stages to learn classification and
semantic reasoning jointly.

\begin{figure}[ht]
\centering
\begin{tikzpicture}[
  blk/.style={rectangle, draw, rounded corners=2pt, text centered,
              text width=2.1cm, minimum height=0.6cm, font=\scriptsize,
              inner sep=2pt},
  arr/.style={thick, -{Stealth}},
  node distance=0.35cm
]
% Left column: Annotation Phase
\node[blk, fill=blue!12]  (raw)    {ViTHSD\\(10K comments)};
\node[blk, fill=orange!20, below=of raw]   (gem)    {Gemini 2.0 Flash\\CoT Prompting};
\node[blk, fill=orange!20, below=of gem]   (filt)   {Quality\\Filtering};
\node[blk, fill=green!20,  below=of filt]  (rdata)  {Rationale Set\\(1,221 pairs)};

% Right column: Training Phase
\node[blk, fill=purple!12, right=2.2cm of raw]   (s1) {Stage 1\\Classifier\\(7,540 data)};
\node[blk, fill=purple!22, below=of s1]           (s2) {Stage 2\\Semantic Align.\\(1,221 data)};
\node[blk, fill=red!12,    below=of s2]           (mdl){HARE Model\\Qwen2.5-3B};

% Arrows within each column
\draw[arr] (raw)   -- (gem);
\draw[arr] (gem)   -- (filt);
\draw[arr] (filt)  -- (rdata);
\draw[arr] (s1)    -- (s2);
\draw[arr] (s2)    -- (mdl);

% Cross-column arrows
\draw[arr] (raw.east)   to[bend left=18]  node[above, font=\tiny, xshift=4pt]{labels} (s1.west);
\draw[arr] (rdata.east) to[bend right=18] (s2.west);

% Phase labels
\node[font=\tiny\itshape, above=2pt of raw]  {\textbf{Annotation Phase}};
\node[font=\tiny\itshape, above=2pt of s1]   {\textbf{Training Phase}};
\end{tikzpicture}
\caption{HARE Framework: Annotation pipeline (left) and two-stage training (right).}
\label{fig:pipeline}
\end{figure}

\subsection{Prompt Engineering for Rationale Generation}

Manually annotating thousands of implicit hate comments with rationales is
prohibitively expensive. We therefore employ Gemini 2.0 Flash \cite{google2025gemini2flash}
as an automated annotator. Gemini is selected for its strong multilingual capabilities,
instruction adherence, and low inference cost at the required scale. We developed a
rigorous three-phase prompt engineering workflow to ensure annotation quality.

\textbf{Phase 1 (Base Prompt):} A simple Role-Task-Output structure. The model
paraphrased inputs rather than inferring hidden meaning; rationales were too long,
unfocused, and in an inconsistent format. Result: rejected.

\textbf{Phase 2 (Psychologist Role):} Added few-shot examples and instructed the model
to adopt a psychologist's perspective. This caused domain drift---the model over-analyzed
benign anger as hate speech (false positives) and wasted tokens on unnecessary
grammatical analysis. Result: rejected.

\textbf{Phase 3 (Structured CoT---Final):} The officially used version. A strict
four-step reasoning structure is enforced: (1) \textit{Target}---identify the target
entity; (2) \textit{Implied}---describe the hidden meaning (3--8 words maximum);
(3) \textit{Evidence}---extract specific offensive words or phrases;
(4) \textit{Verdict}---assign the label based on accumulated evidence. Output is
enforced as structured JSON. This version produced the highest-quality annotations.

\subsection{Rationale Dataset Construction and Filtering}

LLM-generated annotations are not guaranteed to be noise-free. We apply a two-step
filtering process to ensure data quality:

\begin{enumerate}[leftmargin=*]
\item \textbf{Manual Verification:} A stratified random sample is reviewed by all 5 team
members to assess implied statement quality, yielding a final round of prompt
refinement before full generation.
\item \textbf{Automatic Consistency Check:} Each (comment, generated rationale) pair is
re-fed into Gemini Flash with the standard ViTHSD classification prompt. If the
re-predicted label matches the ground truth, the rationale is logically sound and
retained; otherwise it is discarded. This self-consistency check distilled the generated
corpus to \textbf{1,221 high-quality (comment, implied statement, rationale, label)
tuples}.
\end{enumerate}

\subsection{Model Selection and Justification}

We select \textbf{Qwen2.5-3B-Instruct} \cite{Yang2025Qwen25} as our backbone for three
reasons. First, at 3B parameters it fits on a single NVIDIA T4 GPU (16\,GB) under 4-bit
quantization, enabling cost-effective experiment iteration on Kaggle and Colab. Second,
it achieves state-of-the-art results on Vietnamese benchmarks among sub-10B models,
outperforming Llama~2-7B \cite{llama} and Gemma-7B \cite{Mesnard2024Gemma} on
Vietnamese tasks. Third, as an instruction-following model it natively accommodates the
structured CoT output format (JSON with Target, Implied, Evidence, Verdict fields)
required by our framework without additional prefix tuning.

\subsection{Two-Stage QLoRA Fine-Tuning Strategy}

We propose separating classification learning from semantic reasoning learning across
two fine-tuning stages, both implemented via QLoRA \cite{Dettmers2023QLoRA} in 4-bit
NF4 quantization with LoRA rank $r=16$, $\alpha=32$, dropout 0.05.

\textbf{Stage 1 (Classifier Training):} The model learns to predict
\texttt{[Target]\#[Level]} from the full ViTHSD training set. To mitigate class
imbalance, we apply targeted oversampling (+540 samples) for Politics and Religion
minority classes, yielding 7,540 training examples. Hyperparameters: 2 epochs, batch
size 4, learning rate $2\times 10^{-4}$.

\textbf{Stage 2 (Semantic Alignment):} Starting from the Stage~1 checkpoint, we
continue fine-tuning on the 1,221-sample Rationale Dataset. The model jointly learns to
generate the implied statement and classification rationale \textit{before} emitting the
final label. This aligns the model's embedding space to associate implicit expressions
with their true semantic meanings---e.g., that ``99\% power'' in a gaming context
implies \textit{weakness}. This stage is the key differentiator of our approach.
Hyperparameters: 3 epochs, same configuration as Stage~1.

% ==========================================================================
\section{Experiments and Results}\label{sec:results}
% ==========================================================================

This section evaluates HARE against four baselines under identical conditions and
analyzes results from quantitative and linguistic perspectives.

\subsection{Experimental Setup}

All models are evaluated on the ViTHSD test set (1,800 samples) in a multi-label
classification setting. We report F1-Micro (frequency-weighted), F1-Macro (class-equal
weighting, more sensitive to rare categories), Precision, and Recall. Each
configuration is trained and evaluated over five independent runs with different random
seeds, and we report mean and standard deviation ($\mu\pm\sigma$). This protocol
reduces seed sensitivity and better reflects model stability on rare labels.

The four baselines, each representing a distinct modeling paradigm, are:

\begin{enumerate}[leftmargin=*]
\item \textbf{PhoBERT-base} \cite{nguyen2020phobert}: The dominant encoder-only model for
Vietnamese. Represents the standard supervised fine-tuning paradigm on formal Vietnamese
text.
\item \textbf{Flan-T5-base} \cite{chung2022scaling}: An instruction-tuned encoder-decoder
model. Represents the text-to-text paradigm with CoT-style prompting.
\item \textbf{Qwen2.5-3B Vanilla} \cite{Yang2025Qwen25}: Our backbone with Stage~1
training only (no rationale data). This ablation isolates the contribution of semantic
alignment.
\item \textbf{HARE (proposed)}: Qwen2.5-3B with both Stage~1 and Stage~2 training,
including rationale and implied statement data.
\end{enumerate}

To better approximate broader comparison expectations while remaining within the
6-page limit, we include a concise \textit{contextual comparison} against recent
Vietnamese hate speech systems reported in ViHSD/ViTHSD studies
\cite{10.1007/978-3-030-79457-6_35,vithsd}, and use our four retrained systems for
strict apples-to-apples evaluation under a unified preprocessing and label space.

\subsection{Quantitative Results}

Table~\ref{tab:overall_perf} presents the overall multi-label performance comparison on
the ViTHSD test set. Under the 5-run protocol, \textbf{HARE achieves the best mean
F1-Micro of 60.26\%}, outperforming
PhoBERT by +6.14\% and Flan-T5 by +13.42\%. Compared to Qwen2.5 Vanilla, Stage~2
semantic alignment improves F1-Micro by +1.26\% and Precision-Micro by +2.47\%. Higher
precision is particularly important in content moderation, where incorrectly flagging
benign comments negatively impacts user experience.

\begin{table}[ht]
\centering
\caption{Overall Multi-Label Performance on ViTHSD Test Set}
\label{tab:overall_perf}
\scriptsize
\setlength{\tabcolsep}{2pt}
\begin{tabular}{lcccc}
\toprule
\textbf{Metric} & \textbf{HARE} & \textbf{Qwen Van.} & \textbf{PhoBERT} & \textbf{Flan-T5} \\
\midrule
F1-Samples  & \textbf{0.6345} & 0.6332          & 0.5801 & 0.5074 \\
F1-Micro    & \textbf{0.6026} & 0.5900          & 0.5412 & 0.4684 \\
F1-Macro    & 0.3035          & \textbf{0.3310} & 0.2586 & 0.1311 \\
\midrule
Prec-Micro  & \textbf{0.6347} & 0.6100          & 0.5620 & 0.4810 \\
Prec-Macro  & \textbf{0.4067} & 0.3800          & 0.3200 & 0.1800 \\
\midrule
Rec-Micro   & \textbf{0.5735} & 0.5600          & 0.5310 & 0.4520 \\
Rec-Macro   & \textbf{0.2842} & 0.2700          & 0.2450 & 0.1100 \\
\bottomrule
\end{tabular}
\end{table}

\subsection{Per-Label F1 Analysis}

Table~\ref{tab:label_f1} breaks down F1-scores by label category. This per-label
analysis is essential for understanding the linguistic impact of rationale-enhanced
training, particularly on low-resource and culturally nuanced categories.

\begin{table}[ht]
\centering
\caption{Per-Label F1-Score on ViTHSD Test Set (\%)}
\label{tab:label_f1}
\scriptsize
\setlength{\tabcolsep}{2pt}
\begin{tabular}{lcccc}
\toprule
\textbf{Label} & \textbf{HARE} & \textbf{Qwen Van.} & \textbf{PhoBERT} & \textbf{Flan-T5} \\
\midrule
Normal               & \textbf{79.84} & 79.70 & 75.31 & 67.32 \\
Individuals\#Hate    & 61.66 & \textbf{61.90} & 53.65 & 35.83 \\
Groups\#Hate         & \textbf{53.44} & 45.83 & 46.03 & 34.93 \\
Politics\#Hate       & \textbf{54.00} & 36.59 & 29.13 & 2.99  \\
Race\#Hate           & 28.17 & \textbf{32.88} & 28.57 & 3.12  \\
\bottomrule
\end{tabular}
\end{table}

\subsection{Ablation Study: Effect of Two-Stage Training}

To isolate the contribution of each training stage, Table~\ref{tab:ablation} compares
two configurations. Adding Stage~2 semantic alignment consistently improves both
F1-Micro and Precision-Micro, confirming that rationale-guided fine-tuning reduces
false positives and improves overall classification quality, despite mixed results on
the severely data-sparse Race category.

\begin{table}[ht]
\centering
\caption{Ablation Study: Two-Stage Training Contribution}
\label{tab:ablation}
\scriptsize
\setlength{\tabcolsep}{4pt}
\begin{tabular}{lcc}
\toprule
\textbf{Configuration} & \textbf{F1-Micro} & \textbf{Prec-Micro} \\
\midrule
Qwen2.5 Stage 1 only (Vanilla) & 0.5900 & 0.6100 \\
HARE: Stage 1 + Stage 2        & \textbf{0.6026} & \textbf{0.6347} \\
\midrule
$\Delta$ (HARE vs.\ Vanilla)   & \textbf{+1.26\%} & \textbf{+2.47\%} \\
\bottomrule
\end{tabular}
\end{table}

\subsection{Discussion: Linguistic Perspective on Results}

We analyze the experimental results from three linguistic and cultural angles to explain
why HARE succeeds or fails in specific categories.

\textbf{(1) Breakthrough on Political Hate Speech.}
The Politics\#Hate category shows the most dramatic improvement: HARE achieves 54.00\%
F1, a \textbf{+17.41\% gain over Qwen Vanilla} and a \textbf{+24.87\% gain over
PhoBERT}. This category is linguistically the most challenging because political hate in
Vietnamese relies on culturally loaded slang (``ba que,'' ``bò đỏ''), historical
metaphors, and implicit dehumanization. Keyword-based and encoder-only models lack
the cultural grounding to decode these expressions. Our rationale data explicitly teaches
the model to associate these slang terms with their hateful political intent, enabling
semantic generalization far beyond surface-level lexical matching.

\textbf{(2) Improvement on Groups\#Hate.}
HARE achieves 53.44\% F1, a +7.61\% gain over Qwen Vanilla and +7.41\% over PhoBERT.
This gain reflects improved ability to handle collective hate, where attacks target
communities rather than named individuals and require broader contextual inference.

\textbf{(3) Regression on Race\#Hate.}
Race\#Hate shows a slight regression (28.17\%) below Qwen Vanilla (32.88\%). This
failure is attributable to \textit{data sparsity}: only 502 training samples exist for
Race targets, and generated rationales for this category are few and likely noisy. With
insufficient anchor examples, Stage~2 alignment over-fits or introduces misaligned
reasoning. This confirms a key lesson: \textbf{CoT reasoning enhances performance when
sufficient contextual data exist, but cannot compensate for fundamental data scarcity}.

\textbf{(4) Why PhoBERT Underperforms.}
PhoBERT was pre-trained exclusively on formal Vietnamese (Wikipedia, news), creating a
significant domain gap with social media language. Informal slang, abbreviations, and
code-switching are systematically absent from its vocabulary, severely limiting its
generalization. Qwen2.5's exposure to diverse informal web text gives it a natural
advantage in capturing social media pragmatics.

% ==========================================================================
\section{Deployment Architecture}\label{sec:deployment}
% ==========================================================================

To demonstrate real-world applicability, we developed a full-stack web application
for real-time Vietnamese hate speech monitoring from YouTube, illustrated in
Fig.~\ref{fig:system}. The architecture follows a standard Client-Server model ensuring
separation of interface and processing logic.

\begin{figure}[ht]
\centering
\begin{tikzpicture}[
  blk/.style={rectangle, draw, rounded corners=2pt, text centered,
              text width=2.0cm, minimum height=0.55cm, font=\scriptsize,
              inner sep=2pt},
  arr/.style={thick, -{Stealth}},
  node distance=0.35cm
]
\node[blk, fill=gray!15]   (url)  {User Input\\(YouTube URL)};
\node[blk, fill=cyan!20,   below=of url]  (fe)   {Frontend\\React + Vite};
\node[blk, fill=orange!20, below=of fe]   (be)   {Backend\\FastAPI};
\node[blk, fill=yellow!20, left=2.0cm of be]  (yt)   {YouTube\\API v3};
\node[blk, fill=purple!20, right=2.0cm of be] (inf)  {Qwen2.5-3B\\Inference};
\node[blk, fill=green!20,  below=of be]   (disp) {Label + Keyword\\Highlight Display};

\draw[arr] (url)  -- (fe);
\draw[arr] (fe)   -- (be);
\draw[arr] (be)   -- (yt);
\draw[arr] (yt)   -- node[above, font=\tiny]{comments} (be);
\draw[arr] (be)   -- (inf);
\draw[arr] (inf)  -- node[above, font=\tiny]{labels}   (be);
\draw[arr] (be)   -- (disp);
\draw[arr] (disp) to[bend left=55] (fe);
\end{tikzpicture}
\caption{Real-time hate speech monitoring system architecture.}
\label{fig:system}
\end{figure}

The system maps ViTHSD's granular labels to three user-accessible categories:
\textbf{Neutral} (Normal + Clean), \textbf{Offensive}, and \textbf{Hate}, enabling
non-expert users to interpret results without domain knowledge.

\textbf{Keyword Highlighting Module (Rule-based):}
A supplementary rule-based module maintains a blacklist lexicon of high-frequency hate
terms extracted from the ViTHSD training vocabulary. On the React frontend, comments
are scanned and matching terms are visually highlighted, providing fast visual
explanation alongside AI-predicted labels. This aids moderators in quickly pinpointing
dangerous text spans without reading every word.

\textbf{Limitations of Rule-based Highlighting:}
This mechanism only detects explicit keywords and fails entirely on implicit hate that
avoids offensive language (e.g., \foreignlanguage{vietnamese}{\textit{``Khôn như bạn
quê tôi xích đầy''}}). It also generates false positives when blacklisted words appear
in neutral context. These limitations directly motivate the attention-based highlighting
direction described in Section~\ref{sec:conclusion}.

% ==========================================================================
\section{Conclusion and Future Work}\label{sec:conclusion}
% ==========================================================================

This paper presented HARE, a Vietnamese targeted hate speech detection framework, making
five contributions consistent with the goals stated in Section~\ref{sec:intro}. We (C1)
formalized the multi-label targeted hate speech detection task on ViTHSD with
comprehensive linguistic analysis; (C2) designed a three-phase CoT annotation pipeline
using Gemini 2.0 Flash to generate 1,221 high-quality rationale pairs; (C3) proposed a
two-stage QLoRA fine-tuning strategy separating classification and semantic alignment;
(C4) achieved state-of-the-art F1-Micro of 60.26\% on the ViTHSD test set with a
breakthrough +24.87\% gain on Politics\#Hate over PhoBERT; and (C5) delivered a
full-stack real-time monitoring system with visual highlighting. HARE achieves the best
overall F1-Micro (60.26\%), best Precision-Micro (63.47\%), and most significant
improvements on the most linguistically challenging and culturally sensitive categories.

Two principal limitations constrain the current system. First, the rule-based keyword
highlighting fails on implicit hate speech, providing no semantic grounding. Second,
performance on Race\#Hate remains limited by severe data sparsity (502 training
samples), confirming that Chain-of-Thought reasoning cannot substitute for adequate
training data coverage.

Three future directions follow directly from these limitations. (1) \textbf{Attention-based
Highlighting}: Replace the static blacklist with attention weight extraction from
Qwen2.5's self-attention layers, enabling model-guided context-sensitive span
identification including metaphorical expressions. (2) \textbf{Targeted Data
Augmentation}: Apply generative augmentation for underrepresented categories (Race,
Religion) to resolve data sparsity and improve generalization. (3) \textbf{Multimodal
Extension}: Extend the system to handle image-text pairs (memes, thumbnails) for more
comprehensive content moderation on platforms where visual hate speech is prevalent.

% ==========================================================================
\begin{thebibliography}{99}

\bibitem{10.1007/978-3-030-79457-6_35}
S.~T. Luu, K.~V. Nguyen, and N.~L.-T. Nguyen, ``A large-scale dataset for hate speech
detection on Vietnamese social media texts,'' in \textit{Proc. Int. Conf. Intelligent
Information and Database Systems (IEA/AIE)}, Springer, 2021, pp. 415--426.

\bibitem{vithsd}
C.~N. Vo, K.~B. Huynh, S.~T. Luu, and D. Trong-Hop, ``ViTHSD: Exploiting hatred by
targets for hate speech detection on Vietnamese social media texts,'' \textit{J.
Comput. Social Sci.}, 2025, doi: 10.1007/s42001-024-00348-6.

\bibitem{nguyen2020phobert}
D.~Q. Nguyen and A.~T. Nguyen, ``PhoBERT: Pre-trained language models for
Vietnamese,'' in \textit{Findings ACL: EMNLP 2020}, 2020, pp. 1037--1042.

\bibitem{Yang2025Qwen25}
A.~Yang \textit{et al.}, ``Qwen2.5 technical report,'' \textit{arXiv:2412.15115},
2025.

\bibitem{chung2022scaling}
H.~W. Chung \textit{et al.}, ``Scaling instruction-finetuned language models,''
\textit{arXiv:2210.11416}, 2022.

\bibitem{google2025gemini2flash}
Google DeepMind, ``Gemini 2.0 Flash: Efficient multimodal reasoning model,'' Google
AI Studio Documentation, 2025.

\bibitem{llama}
H.~Touvron \textit{et al.}, ``Llama 2: Open foundation and fine-tuned chat models,''
\textit{arXiv:2307.09288}, 2023.

\bibitem{Mesnard2024Gemma}
T.~Mesnard \textit{et al.}, ``Gemma: Open models based on Gemini research and
technology,'' \textit{arXiv:2403.08295}, 2024.

\bibitem{Mikolov2013Efficient}
T.~Mikolov, K.~Chen, G.~Corrado, and J.~Dean, ``Efficient estimation of word
representations in vector space,'' in \textit{Proc. ICLR 2013 Workshop}, 2013.

\bibitem{Bojanowski2017Enriching}
P.~Bojanowski, E.~Grave, A.~Joulin, and T.~Mikolov, ``Enriching word vectors with
subword information,'' \textit{Trans. Assoc. Comput. Linguistics}, vol. 5,
pp. 135--146, 2017.

\bibitem{Yang2023HARE}
Y.~Yang \textit{et al.}, ``HARE: Explainable hate speech detection with step-by-step
reasoning,'' in \textit{Findings ACL: EMNLP 2023}, 2023, pp. 5490--5505.

\bibitem{Wei2022ChainOfThought}
J.~Wei \textit{et al.}, ``Chain-of-thought prompting elicits reasoning in large
language models,'' in \textit{Proc. NeurIPS 2022}, 2022.

\bibitem{Dettmers2023QLoRA}
T.~Dettmers, A.~Pagnoni, A.~Holtzman, and L.~Zettlemoyer, ``QLoRA: Efficient
finetuning of quantized LLMs,'' in \textit{Proc. NeurIPS 2023}, 2023.

\bibitem{Pontiki2014SemEval4}
M.~Pontiki \textit{et al.}, ``SemEval-2014 task 4: Aspect based sentiment
analysis,'' in \textit{Proc. 8th Int. Workshop Semantic Evaluation (SemEval 2014)},
2014, pp. 27--35.

\bibitem{Vaswani2017Attention}
A.~Vaswani \textit{et al.}, ``Attention is all you need,'' in \textit{Proc. NeurIPS
2017}, 2017, pp. 6000--6010.

\bibitem{Devlin2019BERT}
J.~Devlin, M.-W. Chang, K.~Lee, and K.~Toutanova, ``BERT: Pre-training of deep
bidirectional transformers for language understanding,'' in \textit{Proc. NAACL-HLT
2019}, 2019, pp. 4171--4186.

\end{thebibliography}

\end{document}
